%
\againframe{plan}
%
\begin{frame}{Lecture 3 plan}
	\begin{learningobjectives}
		\begin{itemize}
			\item Classification of semiconductors
			\item Nature of bonding in semiconductor
			\item Doping of semiconductor
		\end{itemize}
	\end{learningobjectives}
\end{frame}
%
\begin{frame}{Semiconductors: Classification}
	\begin{definition}
		Semiconductor is a class of condensed matter having electrical conductivity  (\SIrange{e-4}{1}{\siemens\per\metre}) greater than that of insulator \( ( \SI{<e-10}{\siemens\per\metre}) \) but lower than that of conductor (\SI{e7}{\siemens\per\metre}).
	\end{definition}
	%
\begin{columns}
	\begin{column}{0.8\textwidth}
			\begin{itemize}
		\item In the periodic table, there are 11 elements that are semiconductors. These are called \textbf{elemental semiconductors}.  \ce{Si} and \ce{Ge} from Group IV are widely used elemental semiconductors.
	\end{itemize}
	\end{column}
	%
	\begin{column}{0.18\textwidth}
		\begin{figure}
			\includegraphics[width=\textwidth]{semiconductors_periodic_table}
		\end{figure}
	\end{column}
\end{columns}
%
\begin{itemize}
	\item Certain combinations of elements of Groups III and V or Groups II and VI also are semiconductors. These are called \textbf{compound semiconductors}. \ce{GaAs}, \ce{InP} are some of the well known examples.
	%
	\item 	In a semiconductor, two charge carriers, namely \textbf{electrons} and \textbf{holes}, transport current.
\end{itemize}
\end{frame}
%
\begin{frame}{Nature of bonding: Silicon crystal}
\begin{figure}
	\includegraphics[width=0.6\textwidth]{semiconductors_silicon_crystal}
	\caption{(a) Tetrahedral arrangement (b) Diamond crystal structure. Taken from Avadhanulu, Chapter 37.}
\end{figure}
%
\begin{itemize}
	\item \ce{^18Si} atom has electronic configuration \ce{1s^2 2s^2 2p^6 3s^2 3p^2}.
	\item In \ce{Si} crystal, the four valence electrons undergo \textbf{hybridization} to form four \ce{sp^3} molecular orbitals. These orbitals form covalent bonds with neighbouring atoms arranged in tetrahedral arrangement. The resulting structure is the \textbf{diamond cubic crystal structure}.
\end{itemize}
\end{frame}
%
\begin{frame}{Intrinsic Semiconductor}
	\begin{itemize}
		\item Chemically pure semiconductor is called intrinsic semiconductor. \item Practically, a semiconductor is considered pure if there is less than one impurity atom per billion host atoms i.e. $< 1$~ ppb (part per billion).
		\item It has five properties
		%
		\begin{enumerate}
			\item Perfect insulator at absolute zero temperature
			\item Generation of charge carriers by thermalization
			\item Existence of energy band gap (M2U3)
			\item Conductivity is \textbf{highly} influenced by temperature
			\item Recombination of charge carriers
		\end{enumerate}
		\item For \ce{Si}, the valence electron density is of the order of \SI{e29}{\per\metre\cubed}. However, at room temperature the conduction electron density is only of the order of \SI{e18}{\per\metre\cubed}.
	\end{itemize}
	%
	\begin{guesstimate}{Valence electron density of \ce{Si}}
		The lattice constant of Si crystal is \SI{5.43}{\angstrom}.
	\end{guesstimate}
\end{frame}
%
\ifshowhidden
%
\begin{frame}{Nature of metallic bonding}
	\begin{figure}
		\centering
		\includegraphics[width=0.8\textwidth]{semiconductors_sodium_vapour_to_solid}
	\end{figure}
	%
		\begin{itemize}
		\item When the sodium vapour condenses, the valence electrons ``hop'' onto neighbouring lattice sites. One site A looses an electron while the neighbouring site B gains an extra electron.
		\item The reduction in energy due to loss of electron at A is more than the increase in energy due to gain of electron at B. As a result, hopping is preferred resulting in \textbf{metallic bonding}.
		\item Since hopping is non directional, the metallic bond is non-directional.
		\item Thus the formation of metallic bond results in lowering of energy of \ce{Na} solid.
	\end{itemize}
\end{frame}
%
\fi
%
\begin{frame}{Nature of intrinsic semiconductor at \(T=0\) (property 1)}
	\begin{figure}
		\centering
		\includegraphics[width=\textwidth]{semiconductors_silicon_vapour_to_solid}
	\end{figure}
	\begin{itemize}
		\item Valence electrons are strongly attracted to the nucleus. This prevents the formation of free electron gas.
		\item Instead the valence electrons are shared between neighbouring atoms resulting in \textbf{covalent bond}.
		\item As all the valence electrons are involved in bonding, \fcolorbox{red}{white}{there are no ``free'' electrons at absolute zero temperature.}
	\end{itemize}
	%
	\begin{keyinsight}
		Intrinsic semiconductor is a perfect insulator at absolute zero!
	\end{keyinsight}
\end{frame}
%
\begin{frame}{Nature of intrinsic semiconductor at \(T\neq 0\) (property 2)}
	\begin{figure}
		\centering
		\includegraphics[width=0.8\textwidth]{semiconductors_silicon_vapour_to_solid_finite_temperature}
	\end{figure}
	%
	\vspace{-1em}
	%
	\begin{itemize}
{\small
			\item However, at finite temperature, due to thermalization, few bonds are broken, releasing the electrons to the \textbf{free electron sea}.
		\item In the process of bond breaking, the crystal lattice deficient with one electron can be considered as a new type of charge carrier. This is defined as the \textbf{hole}.
		\item The holes are released into the \textbf{free hole sea}.
		\item Every free electron is associated with a free hole. Hence they are called\textbf{ electron-hole pair}.}
	\end{itemize}
	%
	\begin{definition}
		An intrinsic semiconductor is a semiconductor crystal in which the electrical conduction arises due to \textbf{thermally excited} electrons- hole pairs.
	\end{definition}
\end{frame}
%
\begin{frame}{Nature of hole: Drift motion (property 2)}
	%
	\begin{figure}
		\includegraphics[height=0.3\textheight]{semiconductors_holes_analogy}
		\includegraphics[height=0.3\textheight]{semiconductors_hole_drift}
		\caption{(a) Analogy for a hole: Bubbles (b) Drift motion of holes}
	\end{figure}
	%
	\vspace{-1em}
	\begin{itemize}
{\small
			\item A hole is a \textbf{quantum vacancy} that is described as an absence of electron. 
		\item Analogy is the bubbles appearing in carbonated water. The bubbles are absence of water.
		\item Under application of electric field, the bonded electrons move in the opposite direction. The hole appears to move along the electric field.}
	\end{itemize}
	%
	\begin{keyinsight}
		The motion of hole is along the electric field.
	\end{keyinsight}
	%
\end{frame}
%
\begin{frame}{Energy band gap: Chemistry view (property 3)}
	%
	\begin{figure}
		\includegraphics[height=0.25\textwidth]{semiconductors_energy_bang_gap}
	\end{figure}
	\begin{itemize}
{\small
			\item The electron of the covalent bond needs to overcome bonding strength to become free.
		\item If all the bonds are broken, then the crystal vapourizes. The energy needed to vapourize a crystal is called \textbf{cohesive energy}.
		\item The actual energy to free electron is \SI{1.12}{\electron\volt} and is called the \textbf{energy band gap} (M2U3).}
	\end{itemize}
	%
	\begin{guesstimate}{Bonding strength}
		The cohesive energy of Si crystal is \SI{450}{\kilo\joule\per\mole}. Estimate the \ce{Si-Si} bond strength.
	\end{guesstimate}
	%
\end{frame}
%
\begin{frame}{Charge carriers vs temperature (property 4)}
	\begin{figure}
		\includegraphics[height=0.25\textheight]{semiconductors_thermal_energy_vs_band_gap}
	\end{figure}
	\begin{itemize}
{\small
			\item The scale of thermal energy is given by \(k_B T\). At room temperature, the thermal energy of electron is of the order of \SI{26}{\milli\electron\volt}. \hfill [Estimate]
		\item The bond strength can be considered an \textbf{activation barrier} for the generation of electron-hole pair. This is analogous to rate of chemical reaction
		\[
		\text{covalent bond} + \text{thermal energy} \rightarrow \text{electron} + \text{hole}
		\]
		\item Electron-hole pair density increases exponentially with temperature. This follows the \textbf{Arrhenius law}.
		\[
		n \propto \exp(-\frac{1}{k_B T}) \quad p \propto \exp(-\frac{1}{k_B T})
		\]
		}
	\end{itemize}
\end{frame}
%
\begin{frame}{Conductivity vs temperature (property 4)}
	%
	\begin{figure}
		\includegraphics[height=0.4\textheight]{semiconductors_conductivity_vs_temperature}
	\end{figure}
	\begin{itemize}
		\item Since conductivity is proportional to number density, conductivity also follows Arrhenius law.
		\[
		\sigma \propto n \quad \sigma \propto p
		\]
	\end{itemize}
	\begin{keyinsight}
		Conductivity exponentially depends on temperature.
	\end{keyinsight}
\end{frame}
%
\begin{frame}{Recombination of carriers (property 5)}

		\begin{itemize}
			{\footnotesize
		\item We have seen thermal energy to generate electron-hole pairs. This is an example of \textbf{equilibrium} generation
		\item Light can also be used to generate electron-hole pairs. 
		\[
		\text{covalent bond} + \text{light} \rightarrow \text{electron} + \text{hole}
		\]
		However, this is a \textbf{non-equilibrium} generation process.
		\item The reverse process where the electron-hole pairs combine to become part of the covalent lattice is called \textbf{recombination}.
		\[
		\text{electron} + \text{hole} \rightarrow \text{covalent bond} +  \text{light} / \text{lattice vibration}
		\]
		The additional energy is released either as light (\textbf{photon}) or lattice vibration (\textbf{phonon}).
		\item The generation-recombination of charge carriers is a \textbf{dynamic} two way process.
	}
	\end{itemize}

	%
	\begin{keyinsight}
		For analysis, a semiconductor is a collection of electron, holes, photons and phonons and the interaction between them.
	\end{keyinsight}
\end{frame}
%
\begin{frame}{Applications of intrinsic semiconductor: Switch}
\begin{columns}
	\begin{column}[t]{0.6\textwidth}
			\begin{figure}
		\includegraphics[width=\textwidth]{semiconductors_LDR}\\
		\includegraphics[width=0.25\textwidth]{semiconductors_thermistor_component}
		\includegraphics[width=0.25\textwidth]{semiconductors_photoresistor_component}
	\end{figure}
	\end{column}
	%
	\begin{column}[t]{0.38\textwidth}
{\small
			\begin{itemize}
			\item Thermistor (Temperature dependent resistor). Generation of carriers by temperature.
			\item Photo-resistor: (Light Dependent Resistor (LDR)). Generation of carriers by light.
		\end{itemize}
		%
		\begin{table}
			\centering
			\begin{tabular}{l l}
				\toprule
				App	&	Material	\\
				\midrule
				Thermistor	&	\ce{Fe3O4}\\
				Photoresistor	&	\ce{CdSe}\\
				\bottomrule
			\end{tabular}
		\end{table}}
	\end{column}
\end{columns}
%
\end{frame}
%
\begin{frame}{Application of thermistor: Smartphone}
	%
		\begin{figure}
		\includegraphics[width=0.8\textwidth]{qm_sensors_in_smartphone}
		\caption{Image taken from Techbii.com. \href{https://techbii.com/10-most-important-sensors-mobile-phones-uses/}{Click here for link}}
	\end{figure}
	%
	\begin{itemize}
		\item temperature sensor to monitor chip and battery health
	\end{itemize}
\end{frame}
%
\begin{frame}{Extrinsic semiconductor}
	\begin{itemize}
		\item The electrical properties of intrinsic semiconductor can be altered by adding impurities.
		\item The technique of adding controlled impurity is called \textbf{doping} and the impurity added is called \textbf{dopant}.
		\item Doped semiconductor is called \textbf{extrinsic} semiconductor.
		\item Dopant atom substitutes the position of parent atom. The dopants are of two types -- donor and acceptor.

		\item For \ce{Si}, donors are \ce{P}, \ce{As}, etc and acceptors are \ce{B}, \ce{Al}, etc.
		\item For \ce{P} as a donor dopant, four valence electrons contribute to covalent bonding. The fifth electron is loosely bound to the phosphorus atom
		\item For \ce{B} as an acceptor dopant, three valence electrons contribute to covalent bonding.
	\end{itemize}
	%
	\begin{guesstimate}{Dopant concentration}
		\ce{P} doping is 1 ppm. \hfill \(n_{valence} \sim \) \SI{2e23}{\per\centi\metre\cubed}
	\end{guesstimate}
\end{frame}
%
\begin{frame}{Ionization energy of dopants}
	%
	\begin{definition}
		The energy required to ionize a dopant atom is called ionization energy.
	\end{definition}
	%
	{\footnotesize
	\begin{itemize}
		\item Lattice vibrations (phonons) supply the ionization energy (I.E.)
		\item For \ce{P} as donor dopant,
		\[
		\ce{P + I.E. -> P^+ + e}
		\]
		\item Similarly, for \ce{B} as acceptor dopant,
		\[
		\ce{B + I.E. -> B^- + h}
		\]
		\item If thermal energy (\(E_{th}\)) is greater than I.E. \(E_{ion}\), then ions are completely ionized. This is called \textbf{complete} ionization.
		\item For \(E_{th} \ll E_{ion}\), there is no ionization. This is called \textbf{freeze out}.
		\item \textbf{Partial} ionization occurs when \(E_{th} \lesssim E_{ion}\) 
		\item In \ce{Si}, \(E_{ion} \, \sim \) \SI{25}{\milli\electron\volt} and  \(E_{th} \, \sim\)\SI{25}{\milli\electron\volt} at room temperature. Thus, the ions are \textbf{completely ionized}.
	\end{itemize}}
\end{frame}